\documentclass[11pt]{article}
\usepackage{amsmath,amssymb,geometry,enumitem,longtable,array,tcolorbox,xcolor,hyperref}
\geometry{margin=0.75in}
\setlength{\parindent}{0pt}
\setlength{\parskip}{0.55em}

\begin{document}

\title{CAT Quant Verified Practice Batch 01}
\author{Original PYQ-Inspired Practice Set}
\date{}
\maketitle

\section*{1.\ Instructions}
\begin{itemize}
\item This is an original practice batch inspired by CAT Quant PYQ patterns.
\item It is not an official, leaked, or guaranteed CAT paper.
\item Suggested time: 30 to 40 minutes.
\item Calculators are not assumed.
\item Attempt every question on paper before checking the answer key.
\end{itemize}

\section*{2.\ Question Paper}

\subsection*{Question 1}
\textbf{Topic:} Arithmetic \\
\textbf{Subtopic:} Successive Percentage Change \\
\textbf{Difficulty:} Medium-Hard \\
\textbf{Type:} MCQ \\
\textbf{Estimated Time:} 3 minutes \\
\textbf{PYQ-Inspired Archetype:} Restoring original value after multiple percentage swings.

\begin{tcolorbox}
The price of a commodity rose by 25\% in January and then rose again by 20\% in February. In March, the price fell by \(x\%\), and the price at the end of March equalled the price at the start of January. The value of \(x\) is:
\end{tcolorbox}
\begin{enumerate}[label=(\Alph*)]
\item \(33\dfrac{1}{3}\%\)
\item \(40\%\)
\item \(45\%\)
\item \(50\%\)
\end{enumerate}

\subsection*{Question 2}
\textbf{Topic:} Arithmetic \\
\textbf{Subtopic:} Time, Speed and Distance (Boats and Streams) \\
\textbf{Difficulty:} Hard \\
\textbf{Type:} TITA \\
\textbf{Estimated Time:} 4 minutes \\
\textbf{PYQ-Inspired Archetype:} Two-leg journey with mixed direction times.

\begin{tcolorbox}
A boat takes 4 hours to travel 24 km downstream and 6 hours to travel 24 km upstream. What is the speed of the boat in still water, in km/h?
\end{tcolorbox}
\textbf{Type:} TITA

\subsection*{Question 3}
\textbf{Topic:} Algebra \\
\textbf{Subtopic:} Logarithms \\
\textbf{Difficulty:} Medium-Hard \\
\textbf{Type:} MCQ \\
\textbf{Estimated Time:} 3 minutes \\
\textbf{PYQ-Inspired Archetype:} Multi-base log equation collapsing to a single base.

\begin{tcolorbox}
If \(\log_2 x + \log_4 x + \log_{16} x = \dfrac{21}{4}\), then the value of \(x\) is:
\end{tcolorbox}
\begin{enumerate}[label=(\Alph*)]
\item \(4\)
\item \(8\)
\item \(16\)
\item \(32\)
\end{enumerate}

\subsection*{Question 4}
\textbf{Topic:} Number System \\
\textbf{Subtopic:} Remainders (CRT) \\
\textbf{Difficulty:} Hard \\
\textbf{Type:} TITA \\
\textbf{Estimated Time:} 4 minutes \\
\textbf{PYQ-Inspired Archetype:} Smallest integer with three given remainders, no constant alignment.

\begin{tcolorbox}
Find the smallest positive integer \(N\) that leaves remainder 2 when divided by 5, remainder 3 when divided by 7, and remainder 4 when divided by 9.
\end{tcolorbox}
\textbf{Type:} TITA

\subsection*{Question 5}
\textbf{Topic:} Geometry \\
\textbf{Subtopic:} Right Triangle and Inscribed Circle \\
\textbf{Difficulty:} Medium-Hard \\
\textbf{Type:} MCQ \\
\textbf{Estimated Time:} 3 minutes \\
\textbf{PYQ-Inspired Archetype:} Triangle area minus inscribed-circle area.

\begin{tcolorbox}
In a right-angled triangle, the lengths of the two legs are 9 cm and 12 cm. A circle is inscribed in the triangle. What is the area, in sq cm, of the region inside the triangle but outside the inscribed circle?
\end{tcolorbox}
\begin{enumerate}[label=(\Alph*)]
\item \(54 - 9\pi\)
\item \(54 - 7\pi\)
\item \(54 - 6\pi\)
\item \(54 - 4\pi\)
\end{enumerate}

\subsection*{Question 6}
\textbf{Topic:} Arithmetic \\
\textbf{Subtopic:} Mixtures (Repeated Replacement) \\
\textbf{Difficulty:} Medium-Hard \\
\textbf{Type:} MCQ \\
\textbf{Estimated Time:} 3 minutes \\
\textbf{PYQ-Inspired Archetype:} Standard repeated-replacement formula trap.

\begin{tcolorbox}
A vessel contains 80 litres of pure milk. 16 litres of the contents is drawn out and replaced with water; this process is repeated two more times (with 16 litres each time). After the three operations, what percentage of the final mixture is milk?
\end{tcolorbox}
\begin{enumerate}[label=(\Alph*)]
\item \(40.0\%\)
\item \(48.0\%\)
\item \(51.2\%\)
\item \(64.0\%\)
\end{enumerate}

\subsection*{Question 7}
\textbf{Topic:} Algebra \\
\textbf{Subtopic:} Quadratic Equations \\
\textbf{Difficulty:} Hard \\
\textbf{Type:} TITA \\
\textbf{Estimated Time:} 4 minutes \\
\textbf{PYQ-Inspired Archetype:} Integer roots with a Vieta-style ratio constraint.

\begin{tcolorbox}
Let \(p\) and \(q\) be positive integers such that the equation \(x^2 - p x + q = 0\) has two distinct positive integer roots, with one root being exactly twice the other. If \(q = 50\), find the value of \(p\).
\end{tcolorbox}
\textbf{Type:} TITA

\subsection*{Question 8}
\textbf{Topic:} Arithmetic \\
\textbf{Subtopic:} Time and Work (Alternate Days) \\
\textbf{Difficulty:} Medium-Hard \\
\textbf{Type:} MCQ \\
\textbf{Estimated Time:} 3 minutes \\
\textbf{PYQ-Inspired Archetype:} Alternate-day work with exact final-day fit.

\begin{tcolorbox}
A alone can complete a job in 15 days, and B alone can complete the same job in 20 days. They work on alternate days, with A starting on day 1. The job is completed at the end of which day?
\end{tcolorbox}
\begin{enumerate}[label=(\Alph*)]
\item Day 16 \quad
\item Day 17 \quad
\item Day 18 \quad
\item Day 19
\end{enumerate}

\subsection*{Question 9}
\textbf{Topic:} Algebra \\
\textbf{Subtopic:} Inequalities with Modulus \\
\textbf{Difficulty:} Hard \\
\textbf{Type:} MCQ \\
\textbf{Estimated Time:} 4 minutes \\
\textbf{PYQ-Inspired Archetype:} Sum-of-distances inequality on the number line.

\begin{tcolorbox}
The number of integer values of \(x\) satisfying \(|x-2| + |x+3| \le 9\) is:
\end{tcolorbox}
\begin{enumerate}[label=(\Alph*)]
\item 9
\item 10
\item 11
\item 12
\end{enumerate}

\subsection*{Question 10}
\textbf{Topic:} Number System \\
\textbf{Subtopic:} Digits and Divisibility \\
\textbf{Difficulty:} Hard \\
\textbf{Type:} TITA \\
\textbf{Estimated Time:} 4 minutes \\
\textbf{PYQ-Inspired Archetype:} Count of integers with digit-sum and distinct-digit constraints.

\begin{tcolorbox}
How many three-digit positive integers are divisible by 9 and have all three digits distinct and non-zero?
\end{tcolorbox}
\textbf{Type:} TITA

\subsection*{Question 11}
\textbf{Topic:} Geometry \\
\textbf{Subtopic:} Coordinate Geometry (Reflection) \\
\textbf{Difficulty:} Medium-Hard \\
\textbf{Type:} MCQ \\
\textbf{Estimated Time:} 3 minutes \\
\textbf{PYQ-Inspired Archetype:} Reflection of a point about a slanted line.

\begin{tcolorbox}
The reflection of the point \((4, 1)\) about the line \(x + y = 7\) is:
\end{tcolorbox}
\begin{enumerate}[label=(\Alph*)]
\item \((6, 3)\)
\item \((3, 6)\)
\item \((5, 4)\)
\item \((4, 5)\)
\end{enumerate}

\subsection*{Question 12}
\textbf{Topic:} Arithmetic \\
\textbf{Subtopic:} Profit, Loss and Successive Discount \\
\textbf{Difficulty:} Hard \\
\textbf{Type:} MCQ \\
\textbf{Estimated Time:} 4 minutes \\
\textbf{PYQ-Inspired Archetype:} Markup followed by two successive discounts.

\begin{tcolorbox}
A shopkeeper marks up an article by 60\% above cost, offers a 10\% discount on the marked price, and then offers a further 5\% discount on the already discounted price. The overall profit percentage on cost is:
\end{tcolorbox}
\begin{enumerate}[label=(\Alph*)]
\item 32.0\%
\item 36.8\%
\item 38.0\%
\item 45.0\%
\end{enumerate}

\subsection*{Question 13}
\textbf{Topic:} Algebra \\
\textbf{Subtopic:} Arithmetic Progression and Inclusion-Exclusion \\
\textbf{Difficulty:} Hard \\
\textbf{Type:} MCQ \\
\textbf{Estimated Time:} 4 minutes \\
\textbf{PYQ-Inspired Archetype:} Sum of integers in a range excluding two divisibility classes.

\begin{tcolorbox}
The sum of all positive integers from 1 to 200 (inclusive) that are neither divisible by 3 nor divisible by 5 is:
\end{tcolorbox}
\begin{enumerate}[label=(\Alph*)]
\item 10632
\item 10732
\item 10832
\item 10932
\end{enumerate}

\subsection*{Question 14}
\textbf{Topic:} Geometry \\
\textbf{Subtopic:} Mensuration (Volume Conservation) \\
\textbf{Difficulty:} Medium-Hard \\
\textbf{Type:} MCQ \\
\textbf{Estimated Time:} 3 minutes \\
\textbf{PYQ-Inspired Archetype:} Melting a cylinder into spheres.

\begin{tcolorbox}
A solid right circular cylinder of radius 3 cm and height 8 cm is melted and recast into solid spheres each of radius 1 cm. Assuming no material loss, the number of spheres formed is:
\end{tcolorbox}
\begin{enumerate}[label=(\Alph*)]
\item 24
\item 36
\item 48
\item 54
\end{enumerate}

\subsection*{Question 15}
\textbf{Topic:} Arithmetic \\
\textbf{Subtopic:} Averages (Membership Change) \\
\textbf{Difficulty:} Hard \\
\textbf{Type:} MCQ \\
\textbf{Estimated Time:} 4 minutes \\
\textbf{PYQ-Inspired Archetype:} Average shift when one member leaves and two join.

\begin{tcolorbox}
The average weight of a group of 10 people is 62 kg. One person weighing 75 kg leaves the group, and two new people whose average weight is \(w\) kg join the group. The average weight of the resulting group of 11 people becomes 63 kg. The value of \(w\) (in kg) is:
\end{tcolorbox}
\begin{enumerate}[label=(\Alph*)]
\item 70
\item 72
\item 74
\item 76
\end{enumerate}

\newpage
\section*{3.\ Answer Key}

\begin{center}
\begin{tabular}{|c|l|c|c|c|}
\hline
\textbf{Q.No.} & \textbf{Topic} & \textbf{Type} & \textbf{Answer} & \textbf{Difficulty} \\
\hline
1  & Arithmetic - Percentages             & MCQ  & (A) \(33\tfrac{1}{3}\%\)  & Medium-Hard \\ \hline
2  & Arithmetic - TSD (Boats)             & TITA & \(5\) km/h                & Hard \\ \hline
3  & Algebra - Logarithms                 & MCQ  & (B) \(8\)                 & Medium-Hard \\ \hline
4  & Number System - Remainders           & TITA & \(157\)                   & Hard \\ \hline
5  & Geometry - Right Triangle            & MCQ  & (A) \(54-9\pi\)           & Medium-Hard \\ \hline
6  & Arithmetic - Mixtures                & MCQ  & (C) \(51.2\%\)            & Medium-Hard \\ \hline
7  & Algebra - Quadratic Equations        & TITA & \(15\)                    & Hard \\ \hline
8  & Arithmetic - Time and Work           & MCQ  & (B) Day 17                & Medium-Hard \\ \hline
9  & Algebra - Modulus Inequality         & MCQ  & (B) 10                    & Hard \\ \hline
10 & Number System - Digits               & TITA & \(60\)                    & Hard \\ \hline
11 & Geometry - Coordinate (Reflection)   & MCQ  & (A) \((6,3)\)             & Medium-Hard \\ \hline
12 & Arithmetic - Profit and Loss         & MCQ  & (B) \(36.8\%\)            & Hard \\ \hline
13 & Algebra - AP/Inclusion-Exclusion     & MCQ  & (B) \(10732\)             & Hard \\ \hline
14 & Geometry - Mensuration               & MCQ  & (D) \(54\)                & Medium-Hard \\ \hline
15 & Arithmetic - Averages                & MCQ  & (C) \(74\) kg             & Hard \\ \hline
\end{tabular}
\end{center}

\newpage
\section*{4.\ Detailed Solutions}

\subsection*{Solution 1}
\textbf{Answer:} (A) \(33\tfrac{1}{3}\%\)

\textbf{Detailed Solution:}

Step 1: Let the starting price be \(100\).

Step 2: After January (+25\%): \(100 \times 1.25 = 125\).

Step 3: After February (+20\%): \(125 \times 1.20 = 150\).

Step 4: In March the price falls by \(x\%\), and equals \(100\):
\[
150 \times \left(1 - \frac{x}{100}\right) = 100 \;\Rightarrow\; 1 - \frac{x}{100} = \frac{2}{3} \;\Rightarrow\; \frac{x}{100} = \frac{1}{3}.
\]

Step 5: So \(x = \dfrac{100}{3} = 33\tfrac{1}{3}\%\).

\textbf{Why this is CAT-level:}
The trap is to add the percentage changes (\(+25 + 20 = 45\)) and write \(x = 45\%\). Successive percentages must be multiplied.

\textbf{Common Wrong Approach:}
Treating \(\%\) changes as additive on the original base of 100.

\subsection*{Solution 2}
\textbf{Answer:} 5 km/h

\textbf{Detailed Solution:}

Step 1: Downstream speed \(= 24/4 = 6\) km/h. Upstream speed \(= 24/6 = 4\) km/h.

Step 2: Speed in still water \(= \dfrac{\text{downstream} + \text{upstream}}{2} = \dfrac{6+4}{2} = 5\) km/h.

Step 3 (verify): Stream speed \(= (6-4)/2 = 1\) km/h. Downstream \(= 5+1 = 6\) \checkmark. Upstream \(= 5-1 = 4\) \checkmark.

\textbf{Why this is CAT-level:}
Tests the boats-and-streams identity. The trap is averaging the two travel times instead of the two speeds.

\textbf{Common Wrong Approach:}
Computing average speed as \(\dfrac{24+24}{4+6}=4.8\) km/h, which is harmonic, not still-water speed.

\subsection*{Solution 3}
\textbf{Answer:} (B) \(8\)

\textbf{Detailed Solution:}

Step 1: Let \(a = \log_2 x\). Then \(\log_4 x = a/2\) and \(\log_{16} x = a/4\).

Step 2: Equation becomes \(a + \dfrac{a}{2} + \dfrac{a}{4} = \dfrac{21}{4}\), i.e. \(\dfrac{7a}{4} = \dfrac{21}{4}\Rightarrow a = 3\).

Step 3: So \(\log_2 x = 3 \Rightarrow x = 2^3 = 8\).

\textbf{Why this is CAT-level:}
Requires confident change-of-base reduction to a single variable.

\textbf{Common Wrong Approach:}
Multiplying logs instead of using change of base.

\subsection*{Solution 4}
\textbf{Answer:} \(157\)

\textbf{Detailed Solution:}

Step 1: Note that the remainders 2, 3, 4 are \emph{not} a constant amount less than 5, 7, 9 (they are short by 3, 4, 5), so the standard ``\(N + k\) divisible by all'' shortcut does not apply. Solve by successive substitution (CRT).

Step 2: From \(N \equiv 2 \pmod 5\): \(N = 5a + 2\) for \(a \ge 0\), giving \(2, 7, 12, 17, \ldots\).

Step 3: Impose \(N \equiv 3 \pmod 7\). Among \(5a+2\), seek the smallest with \(5a+2\equiv 3\pmod7\), i.e. \(5a \equiv 1 \pmod 7\). Since \(5 \cdot 3 = 15 \equiv 1\pmod 7\), \(a \equiv 3 \pmod 7\). Smallest \(a = 3\Rightarrow N = 17\). General: \(N = 17 + 35 k\), giving \(17, 52, 87, 122, 157, \ldots\).

Step 4: Impose \(N \equiv 4 \pmod 9\):
\(17 \bmod 9 = 8\); \(52 \bmod 9 = 7\); \(87 \bmod 9 = 6\); \(122 \bmod 9 = 5\); \(157 \bmod 9 = 4\). \checkmark

Step 5: Smallest such \(N\) is \(\boxed{157}\).

Verify: \(157 = 5\cdot 31 + 2 = 7\cdot 22 + 3 = 9\cdot 17 + 4\). \checkmark

\textbf{Why this is CAT-level:}
Two-stage CRT, where the ``\(N + k\) clean'' shortcut fails and one must combine constraints systematically.

\textbf{Common Wrong Approach:}
Computing LCM\((5,7,9) = 315\) and adding a fixed constant, missing that the remainders are not aligned.

\subsection*{Solution 5}
\textbf{Answer:} (A) \(54 - 9\pi\)

\textbf{Detailed Solution:}

Step 1: Hypotenuse \(= \sqrt{9^2 + 12^2} = 15\).

Step 2: Area of triangle \(= \tfrac{1}{2}\cdot 9\cdot 12 = 54\) sq cm.

Step 3: Inradius of a right triangle \(= \dfrac{a + b - c}{2} = \dfrac{9 + 12 - 15}{2} = 3\) cm.

Step 4: Area of inscribed circle \(= \pi r^2 = 9\pi\).

Step 5: Required area \(= 54 - 9\pi\).

\textbf{Why this is CAT-level:}
Uses the right-triangle inradius identity; otherwise one must compute the inradius via \(\text{Area} = r \cdot s\) with \(s = 18\), which also gives \(r = 3\).

\textbf{Common Wrong Approach:}
Using circumradius (\(c/2 = 7.5\)) in place of inradius.

\subsection*{Solution 6}
\textbf{Answer:} (C) \(51.2\%\)

\textbf{Detailed Solution:}

Step 1: Fraction of mixture removed each step \(= 16/80 = 1/5\). So remaining fraction \(= 4/5\).

Step 2: After 3 replacements, fraction of original milk remaining \(= (4/5)^3 = 64/125\).

Step 3: Percentage \(= 64/125 \times 100 = 51.2\%\).

\textbf{Why this is CAT-level:}
Repeated-replacement formula. The trap is to subtract \(3 \times 20\% = 60\%\) and write \(40\%\).

\textbf{Common Wrong Approach:}
Linear-subtraction misuse; or computing \((1/5)^3\) for milk fraction.

\subsection*{Solution 7}
\textbf{Answer:} \(15\)

\textbf{Detailed Solution:}

Step 1: Let the roots be \(\alpha\) and \(\beta\), distinct positive integers with \(\alpha = 2\beta\), \(\alpha > \beta\).

Step 2: Product \(\alpha\beta = q = 50\Rightarrow 2\beta^2 = 50\Rightarrow \beta^2 = 25\Rightarrow \beta = 5\) (positive).

Step 3: So \(\alpha = 10\), \(\beta = 5\). Sum \(\alpha + \beta = 15 = p\).

Verify: \(x^2 - 15x + 50 = (x-5)(x-10)\). \checkmark

\textbf{Why this is CAT-level:}
Combines Vieta's relations with a ratio constraint and an integer-existence check.

\textbf{Common Wrong Approach:}
Trying every factor pair of 50 (\((50,1), (25,2), (10,5)\)) without using the doubling condition.

\subsection*{Solution 8}
\textbf{Answer:} (B) Day 17

\textbf{Detailed Solution:}

Step 1: A's daily rate \(= 1/15\). B's daily rate \(= 1/20\).

Step 2: In every 2-day cycle (A on odd day, B on even day), work done \(= 1/15 + 1/20 = 4/60 + 3/60 = 7/60\).

Step 3: After 8 cycles (16 days), total work \(= 8 \times 7/60 = 56/60 = 14/15\).

Step 4: Remaining work \(= 1 - 14/15 = 1/15\). Day 17 is A's turn; A's daily output is exactly \(1/15\).

Step 5: A completes the remaining work in exactly 1 full day. So the job ends at the end of Day 17.

\textbf{Why this is CAT-level:}
The numbers are crafted so the remaining fraction equals exactly one day's output of the day-17 worker — a classic alternating-day trap.

\textbf{Common Wrong Approach:}
Adding rates and dividing 1 by \(7/120\) (treating it as combined-rate), giving \(\approx 17.14\) days.

\subsection*{Solution 9}
\textbf{Answer:} (B) 10

\textbf{Detailed Solution:}

Step 1: \(|x-2|+|x+3|\) is the sum of distances from \(x\) to 2 and to \(-3\).

Step 2: For \(x \in [-3, 2]\), this sum is exactly \(5 \le 9\). All such \(x\) work.

Step 3: For \(x > 2\): expression \(= (x-2) + (x+3) = 2x+1 \le 9 \Rightarrow x \le 4\). So \(x \in (2, 4]\).

Step 4: For \(x < -3\): expression \(= -(x-2) - (x+3) = -2x - 1 \le 9 \Rightarrow x \ge -5\). So \(x \in [-5, -3)\).

Step 5: Combined range: \(x \in [-5, 4]\). Integers: \(-5, -4, -3, -2, -1, 0, 1, 2, 3, 4\), total \(= 10\).

\textbf{Why this is CAT-level:}
The geometric (distance) interpretation collapses three cases into one inequality.

\textbf{Common Wrong Approach:}
Solving \((x-2)+(x+3) \le 9\) for all \(x\) (ignoring the modulus signs in each region).

\subsection*{Solution 10}
\textbf{Answer:} \(60\)

\textbf{Detailed Solution:}

Step 1: Divisibility by 9 requires digit sum to be a multiple of 9. For three distinct non-zero digits chosen from \(\{1,\ldots,9\}\), the digit sum lies between \(1+2+3=6\) and \(7+8+9=24\). The multiples of 9 in this range: 9 and 18.

Step 2: Triples summing to 9 (with distinct elements in \(\{1,\ldots,9\}\)): \(\{1,2,6\}, \{1,3,5\}, \{2,3,4\}\). Count: 3.

Step 3: Triples summing to 18 (with distinct elements in \(\{1,\ldots,9\}\)): enumerate by largest element.

Largest = 9: need pair from \(\{1,\ldots,8\}\) summing to 9 with both \(< 9\): \((1,8),(2,7),(3,6),(4,5)\). Count: 4.

Largest = 8: need pair from \(\{1,\ldots,7\}\) summing to 10: \((3,7),(4,6)\). Count: 2.

Largest = 7: need pair from \(\{1,\ldots,6\}\) summing to 11: \((5,6)\). Count: 1.

Largest \(\le 6\): pair must sum to 12 with both \(< 6\); maximum is \(4+5 = 9\). None.

Total for sum 18: \(4 + 2 + 1 = 7\) triples.

Step 4: Total unordered triples \(= 3 + 7 = 10\).

Step 5: Each unordered triple of three distinct nonzero digits yields \(3! = 6\) three-digit numbers (no zero, so all permutations are valid 3-digit numbers).

Step 6: Total \(= 10 \times 6 = 60\).

\textbf{Why this is CAT-level:}
Combines digit-sum divisibility, restricted digit set, and permutation counting — with the no-zero restriction crucial.

\textbf{Common Wrong Approach:}
Forgetting to exclude triples containing 0 (the problem requires nonzero), or counting ordered triples directly without dividing by symmetries.

\subsection*{Solution 11}
\textbf{Answer:} (A) \((6, 3)\)

\textbf{Detailed Solution:}

Step 1: Reflect \(P=(4,1)\) about line \(L: x+y=7\). Normal direction to \(L\) is \((1,1)\).

Step 2: Parametrize the line through \(P\) in the normal direction: \(P + t(1,1) = (4+t, 1+t)\). Find \(t\) where this hits \(L\): \((4+t)+(1+t) = 7 \Rightarrow 2t + 5 = 7 \Rightarrow t = 1\). Foot of perpendicular \(F = (5, 2)\).

Step 3: Reflection \(P' = 2F - P = (10 - 4,\, 4 - 1) = (6, 3)\).

Verify: midpoint of \(P\) and \(P'\) is \((5,2)\), on \(L\) since \(5+2=7\). \checkmark.\\
Direction \(P' - P = (2, 2)\) is parallel to the normal \((1,1)\). \checkmark.

\textbf{Why this is CAT-level:}
Forces the use of foot-of-perpendicular reflection rather than the simple ``swap coordinates'' rule (which would apply only for \(y = x\)).

\textbf{Common Wrong Approach:}
Reflexively swapping coordinates of \((4,1)\) to give \((1,4)\), ignoring that the mirror is \(x+y=7\), not \(y=x\).

\subsection*{Solution 12}
\textbf{Answer:} (B) \(36.8\%\)

\textbf{Detailed Solution:}

Step 1: Let cost \(C = 100\). Marked price \(= 160\).

Step 2: After 10\% discount: \(160 \times 0.90 = 144\).

Step 3: After further 5\% discount: \(144 \times 0.95 = 136.80\).

Step 4: Profit \(= 136.80 - 100 = 36.80\). Profit \% \(= 36.8\%\).

\textbf{Why this is CAT-level:}
Two successive discounts must be applied multiplicatively, not added.

\textbf{Common Wrong Approach:}
\(60 - 10 - 5 = 45\%\) profit; or \(60 - 15 = 45\%\).

\subsection*{Solution 13}
\textbf{Answer:} (B) \(10732\)

\textbf{Detailed Solution:}

Step 1: Total sum \(S = 1+2+\cdots+200 = \dfrac{200 \cdot 201}{2} = 20100\).

Step 2: Sum of multiples of 3 in \([1,200]\): largest is \(198 = 3 \cdot 66\). Sum \(= 3 \cdot \dfrac{66 \cdot 67}{2} = 3 \cdot 2211 = 6633\).

Step 3: Sum of multiples of 5 in \([1,200]\): largest is \(200 = 5 \cdot 40\). Sum \(= 5 \cdot \dfrac{40 \cdot 41}{2} = 5 \cdot 820 = 4100\).

Step 4: Sum of multiples of 15 in \([1,200]\): largest is \(195 = 15 \cdot 13\). Sum \(= 15 \cdot \dfrac{13 \cdot 14}{2} = 15 \cdot 91 = 1365\).

Step 5: Sum of multiples of 3 or 5 (inclusion-exclusion) \(= 6633 + 4100 - 1365 = 9368\).

Step 6: Required sum \(= 20100 - 9368 = 10732\).

\textbf{Why this is CAT-level:}
Inclusion-exclusion on AP sums; the trap is forgetting to add back the multiples-of-15 overlap.

\textbf{Common Wrong Approach:}
Subtracting both sums without re-adding multiples of 15.

\subsection*{Solution 14}
\textbf{Answer:} (D) \(54\)

\textbf{Detailed Solution:}

Step 1: Volume of cylinder \(= \pi r^2 h = \pi \cdot 9 \cdot 8 = 72\pi\) cubic cm.

Step 2: Volume of one sphere of radius 1 \(= \dfrac{4}{3}\pi(1)^3 = \dfrac{4\pi}{3}\).

Step 3: Number of spheres \(= \dfrac{72\pi}{4\pi/3} = 72 \cdot \dfrac{3}{4} = 54\).

\textbf{Why this is CAT-level:}
Direct volume conservation. The traps are: (a) using diameter as radius, (b) confusing the cone/sphere formula.

\textbf{Common Wrong Approach:}
Using \(\dfrac{4}{3}\pi r^2\) (missing the cube) for sphere volume.

\subsection*{Solution 15}
\textbf{Answer:} (C) \(74\) kg

\textbf{Detailed Solution:}

Step 1: Original total weight \(= 10 \times 62 = 620\) kg.

Step 2: After the 75 kg person leaves: total \(= 620 - 75 = 545\) kg, group size 9.

Step 3: Two new people of average \(w\) join: new total \(= 545 + 2w\), group size \(= 11\).

Step 4: New average \(= 63\): \(\dfrac{545 + 2w}{11} = 63\Rightarrow 545 + 2w = 693\Rightarrow 2w = 148\Rightarrow w = 74\) kg.

Verify: \(545 + 148 = 693\); \(693 / 11 = 63\). \checkmark

\textbf{Why this is CAT-level:}
Tests careful tracking of group-size changes and totals through two operations.

\textbf{Common Wrong Approach:}
Using divisor 10 (forgetting that the group grew from 9 to 11 after two joiners).

\newpage
\section*{5.\ Topic-Wise Distribution}

\begin{center}
\begin{tabular}{|l|c|l|}
\hline
\textbf{Topic} & \textbf{Number of Questions} & \textbf{Question Numbers} \\
\hline
Arithmetic - Percentages         & 1 & 1 \\ \hline
Arithmetic - TSD (Boats)         & 1 & 2 \\ \hline
Arithmetic - Mixtures            & 1 & 6 \\ \hline
Arithmetic - Time and Work       & 1 & 8 \\ \hline
Arithmetic - Profit/Loss         & 1 & 12 \\ \hline
Arithmetic - Averages            & 1 & 15 \\ \hline
Algebra - Logarithms             & 1 & 3 \\ \hline
Algebra - Quadratic Equations    & 1 & 7 \\ \hline
Algebra - Inequalities (Modulus) & 1 & 9 \\ \hline
Algebra - AP / Inclusion-Exclusion & 1 & 13 \\ \hline
Number System - Remainders       & 1 & 4 \\ \hline
Number System - Digits           & 1 & 10 \\ \hline
Geometry - Right Triangle/Circle & 1 & 5 \\ \hline
Geometry - Coordinate            & 1 & 11 \\ \hline
Geometry - Mensuration           & 1 & 14 \\ \hline
\textbf{Total} & \textbf{15} & \\ \hline
\end{tabular}
\end{center}

Topic balance: Arithmetic 6, Algebra 4, Number System 2, Geometry/Mensuration 3.

\section*{6.\ Quality Verification Summary}

\begin{center}
\begin{tabular}{|c|c|c|c|c|}
\hline
\textbf{Question} & \textbf{Answer Verified} & \textbf{MCQ Options Verified} & \textbf{Ambiguity Check} & \textbf{Status} \\
\hline
1  & Yes & Yes            & Clear & Ready \\ \hline
2  & Yes & Not Applicable & Clear & Ready \\ \hline
3  & Yes & Yes            & Clear & Ready \\ \hline
4  & Yes & Not Applicable & Clear & Ready \\ \hline
5  & Yes & Yes            & Clear & Ready \\ \hline
6  & Yes & Yes            & Clear & Ready \\ \hline
7  & Yes & Not Applicable & Clear & Ready \\ \hline
8  & Yes & Yes            & Clear & Ready \\ \hline
9  & Yes & Yes            & Clear & Ready \\ \hline
10 & Yes & Not Applicable & Clear & Ready \\ \hline
11 & Yes & Yes            & Clear & Ready \\ \hline
12 & Yes & Yes            & Clear & Ready \\ \hline
13 & Yes & Yes            & Clear & Ready \\ \hline
14 & Yes & Yes            & Clear & Ready \\ \hline
15 & Yes & Yes            & Clear & Ready \\ \hline
\end{tabular}
\end{center}

\end{document}
